% =====================================================================
%  Árvore datada com o BEAST 2.7
%  Prática 6 -- Disciplina Métodos Computacionais (UNESP)
%
%  Compilar com pdfLaTeX (Overleaf: Menu > Compiler > pdfLaTeX).
%
%  VERSÃO: por padrão gera a versão do ALUNO (sem respostas).
%  Para a versão do PROFESSOR (com gabarito), troque \professorfalse
%  por \professortrue na linha indicada abaixo, ou compile com
%     pdflatex "\def\versaoprofessor{}\input{apostila_beast.tex}"
% =====================================================================
\documentclass[11pt,a4paper]{article}

\newif\ifprofessor
\professorfalse            % <-- troque para \professortrue para incluir o gabarito
\ifdefined\versaoprofessor \professortrue \fi

\usepackage[T1]{fontenc}
\usepackage[utf8]{inputenc}
\usepackage{lmodern}
\usepackage{textcomp}
\usepackage[provide=*,brazilian]{babel}
\usepackage[a4paper,margin=2.5cm]{geometry}
\usepackage{fancyvrb}
\usepackage{upquote}
\usepackage{booktabs}
\usepackage{enumitem}
\usepackage{comment}
\usepackage{url}
\usepackage[hidelinks]{hyperref}

\fvset{xleftmargin=1.5em,fontsize=\small}
\setlist{itemsep=0.2em,topsep=0.4em}
\setlength{\parskip}{0.35em}
\setlength{\emergencystretch}{3em}
\urlstyle{tt}

\newcommand{\arq}[1]{\path{#1}}
\newcommand{\ui}[1]{\textit{#1}}
\newcommand{\tecla}[1]{\textsf{#1}}

\ifprofessor
  \specialcomment{resposta}{\begin{quote}\small\textbf{Gabarito.}\ }{\end{quote}}
\else
  \excludecomment{resposta}
\fi

\title{Árvore datada com o BEAST 2.7\\[0.3em]
\large Prática 6: árvore bayesiana datada com relógio molecular}
\author{Disciplina Métodos Computacionais\\Universidade Estadual Paulista (UNESP)}
\date{\ifprofessor Versão do professor (com gabarito)\else Versão do aluno\fi\\ Novembro de 2026}

\begin{document}
\maketitle

\paragraph{Origem do material.} Esta apostila adapta o tutorial \emph{Building a calibrated gene
tree with BEAST 2.7} (\url{https://labesta.org/blog/2024/calibrated_gene_tree/}) para a matriz
de Hylodidae das práticas anteriores, dividida em duas partições (16S e COI). As mudanças estão
no Apêndice~\ref{ap:mudancas}.

\paragraph{Arquivos.} A pasta \arq{pratica_beast} contém a matriz \arq{hylodidae_concatenado.nex}
(Prática~4), o arquivo \arq{particoes.txt} (Prática~4) e o script \arq{ajuste_taxa.R}.

\paragraph{Um aviso.} Esta é uma análise de tutorial, montada para rodar em poucos minutos.
Ao longo do texto, os quadros ``\textbf{Na vida real}'' indicam o que faríamos diferente numa
análise para publicação.

% =====================================================================
\section{Introdução}
% =====================================================================

O BEAST (\emph{Bayesian Evolutionary Analysis Sampling Trees}) \cite{bouckaert2019} é um
programa de análise filogenética bayesiana que estima árvores com o tempo nos ramos
(cronogramas). Há muitos tutoriais no site do BEAST 2 (\url{https://www.beast2.org}) e no
\emph{Taming the BEAST} (\url{https://taming-the-beast.org}). Uma análise bayesiana combina
três ingredientes: a matriz de dados, o modelo de substituição de nucleotídeos e as
distribuições \emph{a priori} (\emph{priors}) de cada parâmetro. O resultado não é uma única
árvore, mas uma amostra de milhares de árvores e valores de parâmetros, a distribuição
\emph{a posteriori}.

Vamos usar a matriz concatenada de Hylodidae (16S + COI) e, para datar a árvore, uma
\emph{prior} para a taxa de substituição do 16S baseada em estimativas publicadas para anuros,
como no tutorial original.

% =====================================================================
\section{Instalação}
% =====================================================================

\begin{enumerate}
\item \textbf{BEAST 2.} Baixe e instale a versão para o seu sistema em
\url{https://www.beast2.org}. O pacote traz o BEAUti (que monta a análise), o BEAST (que
roda a análise) e o TreeAnnotator (que resume as árvores).
\item \textbf{bModelTest.} Abra o BEAUti, vá em \ui{File > Manage Packages}, selecione
\ui{bModelTest} \cite{bouckaert2017} e clique em \ui{Install/Upgrade}. Feche e abra o BEAUti
de novo para que o pacote seja carregado.
\item \textbf{Tracer.} Baixe em \url{https://github.com/beast-dev/tracer/releases}
\cite{rambaut2018}.
\item \textbf{FigTree.} O mesmo da Prática~5.
\end{enumerate}

% =====================================================================
\section{A matriz com as partições}
% =====================================================================

O BEAUti lê as partições de um bloco \verb|sets| no fim do arquivo NEXUS. Abra o
\arq{hylodidae_concatenado.nex} no Sublime Text ou no Notepad++ e cole no fim do arquivo,
depois do \verb|END;| que fecha a matriz, o bloco abaixo, usando os números do seu
\arq{particoes.txt}:
\begin{Verbatim}
  ;
END;

begin sets;
  charset p16S = 1-2475;
  charset COI  = 2476-3187;
end;
\end{Verbatim}
Salve com o nome \arq{hylodidae_particionado.nex}. O nome \verb|p16S| começa com uma letra
porque, no formato NEXUS, um nome não pode começar com número.

\begin{quote}\small
\textbf{Na vida real.} Na Prática~5, o COI foi dividido nas três posições do códon. Aqui usamos
só duas partições para ter menos parâmetros e uma análise mais rápida.
\end{quote}

% =====================================================================
\section{Montando a análise no BEAUti}
% =====================================================================

\subsection{Partitions}

Vá em \ui{File > Import Alignment} e abra o \arq{hylodidae_particionado.nex}. O BEAUti mostra
uma linha para cada partição: \verb|p16S| (2475 sítios) e \verb|COI| (712).

Os dois fragmentos vêm do DNA mitocondrial, que é herdado como uma unidade só: eles
compartilham a mesma árvore. Selecione as duas linhas (clique na primeira e, com
\tecla{Shift}, na segunda) e clique em \ui{Link Trees}. A coluna \ui{Tree} passa a mostrar
\verb|p16S| nas duas linhas. Os modelos de substituição (\ui{Site Model}) e os relógios
(\ui{Clock Model}) continuam separados: cada marcador terá o seu.

\subsection{Tip Dates}

Deixe como está.

\subsection{Site Model}

Em vez de escolher um modelo de substituição, vamos usar o pacote bModelTest
\cite{bouckaert2017}, que estima o modelo junto com a árvore: durante a análise, ele passa por
vários modelos e calcula a probabilidade de cada um. Na lista da esquerda, clique em \verb|p16S|
e troque \ui{Gamma Site Model} por \ui{BEAST Model Test}. Faça o mesmo para \verb|COI|. Deixe
\ui{Mutation Rate} em 1,0, sem marcar \ui{estimate}: as diferenças de taxa entre os dois
marcadores ficarão por conta dos relógios.

\subsection{Clock Model}\label{sec:clock}

Usaremos o relógio estrito (\ui{Strict Clock}), o padrão do BEAUti, que supõe a mesma taxa de
substituição em todos os ramos da árvore. Cada marcador terá o seu relógio, porque o COI
evolui mais rápido que o 16S.

\begin{quote}\small
\textbf{Na vida real.} Numa árvore com dezenas de milhões de anos, como esta, dificilmente a
taxa é a mesma em todos os ramos. Um relógio relaxado, em que cada ramo tem a sua taxa (por
exemplo, o \ui{Optimised Relaxed Clock}, do pacote ORC), seria provavelmente mais adequado,
mas tem mais parâmetros e demora mais para convergir.
\end{quote}

Para transformar distâncias genéticas em tempo, é preciso informar ao BEAST algo sobre o
tempo: a idade de um nó (por um fóssil ou um evento geológico) ou a taxa de substituição. Como
no tutorial original, vamos usar uma \emph{prior} para a taxa do 16S, ajustada a quatro
estimativas publicadas para anuros \cite{lemmon2007,fouquet2012,gehara2017,dufresnes2022}:
0,006; 0,00555; 0,00277 e 0,005 substituições por sítio por milhão de anos. O script
\arq{ajuste_taxa.R} ajusta uma distribuição lognormal a esses valores:
\begin{Verbatim}
install.packages("fitdistrplus")
library(fitdistrplus)
taxas <- c(0.006, 0.00555, 0.00277, 0.005)
mean(taxas)
ajuste <- fitdist(taxas, "lnorm")
ajuste
\end{Verbatim}
A média das taxas (0,00483) será o valor inicial da taxa do 16S, e os parâmetros da lognormal,
\verb|meanlog| $= -5{,}374295$ e \verb|sdlog| $= 0{,}304072$, serão usados na \emph{prior}
(Seção~\ref{sec:priors}). Como as taxas estão em milhões de anos, as idades da árvore também
sairão em milhões de anos.

O 16S calibra o tempo; a taxa do COI será estimada a partir dos dados. No BEAUti:
\begin{enumerate}
\item No menu \ui{Mode}, desmarque \ui{Automatic set clock rate}.
\item Na aba \ui{Clock Model}, clique em \verb|p16S|, digite \verb|0.00483| em \ui{Mean clock
rate} e marque \ui{estimate}.
\item Clique em \verb|COI|, digite \verb|0.01| em \ui{Mean clock rate} e marque
\ui{estimate}.
\end{enumerate}

\subsection{Priors}\label{sec:priors}

\begin{enumerate}
\item Na primeira linha, \verb|Tree.t:p16S|, deixe o \ui{Yule Model}, a \emph{prior} de árvore
mais simples: as linhagens só se originam, com uma única taxa de especiação. (O \ui{Birth Death
Model} também permite a extinção, mas tem um parâmetro a mais.)

\item \textbf{Taxa do 16S.} Na linha \verb|clockRate.c:p16S|, troque \ui{Uniform} por
\ui{Log Normal} e clique no triângulo à esquerda do nome para abrir os detalhes. Digite
\verb|M| $= -5.374295$ e \verb|S| $= 0.304072$ (sem marcar \ui{Mean In Real Space}). Ao lado,
o BEAUti desenha a distribuição: ela deve ter mediana 0,00463 e média 0,00485. Compare com o
gráfico feito no R pelo fim do script \arq{ajuste_taxa.R}.

\item \textbf{Taxa do COI.} Na linha \verb|clockRate.c:COI|, troque \ui{Uniform} por
\ui{Log Normal} e digite \verb|M| $= -4.68$ e \verb|S| $= 1.0$. Em anfíbios, o COI acumula
diferenças cerca de duas vezes mais depressa que o 16S \cite{vences2005b}. Por isso a mediana
dessa \emph{prior} é o dobro da do 16S ($-5{,}374 + \ln 2 \approx -4{,}68$, mediana 0,00928),
mas o desvio é bem maior ($S = 1$, contra 0,30 do 16S): o intervalo de 95\% vai de 0,0013 a
0,066, ou seja, de cerca de 0,3 a 14 vezes a taxa do 16S. A \emph{prior} do 16S é estreita
porque é ela que calibra o tempo; a do COI é larga porque só ajuda, e quem decide são os dados.

\item \textbf{Monofilia de Hylodidae.} Clique em \ui{+ Add Prior}, escolha \ui{MRCA prior} e
clique em \ui{OK}. Em \ui{Taxon set label}, escreva \verb|Hylodidae|. Selecione todos os táxons
da lista da esquerda (\tecla{Ctrl+A}) e passe-os para a direita com o botão \texttt{>>}; depois, selecione
\emph{Limnomedusa macroglossa} na lista da direita e devolva-a com \texttt{<<}. Clique em \ui{OK}
e, na nova linha (a que começa com o botão \ui{Hylodidae}), marque \ui{monophyletic}.

Temos um único grupo externo, e só um relógio não basta para posicionar bem a raiz da árvore.
Por isso forçamos a monofilia de Hylodidae, que é bem sustentada por estudos anteriores
\cite{vittorazzi2021,desa2022}. Bônus: o Tracer passará a mostrar a idade do ancestral comum
mais recente de Hylodidae (\verb|mrca.age(Hylodidae)|).

\item Deixe as outras \emph{priors} como estão.
\end{enumerate}

\begin{quote}\small
\textbf{Na vida real.} O ideal é não forçar a monofilia e incluir vários grupos externos, de
modo que os dados mostrem se o grupo é monofilético.
\end{quote}

\subsection{MCMC}

Em \ui{Chain Length}, digite \verb|1000000| (um milhão de gerações). Mantenha o registro a cada
1000 gerações (\ui{Log Every}, dentro de \ui{tracelog} e \ui{treelog}), o valor padrão. Salve
o arquivo com \ui{File > Save As}, com o nome \arq{hylodidae.xml}, na pasta \arq{pratica_beast}.

\begin{quote}\small
\textbf{Na vida real.} Um milhão de gerações é pouco: o ideal é rodar pelo menos 10 milhões
($10^7$) e conferir no Tracer se todos os parâmetros convergiram.
\end{quote}

% =====================================================================
\section{Rodando o BEAST}
% =====================================================================

Abra o BEAST, clique em \ui{Choose file...} ao lado de \ui{Input file}, escolha o arquivo
\arq{hylodidae.xml} e clique em \ui{Run}. A janela mostra o andamento da análise, que leva
poucos minutos. Ao terminar, a pasta terá dois arquivos novos:
\begin{itemize}
\item \arq{hylodidae.log}: os valores dos parâmetros a cada 1000 gerações;
\item \arq{hylodidae-p16S.trees}: as árvores amostradas a cada 1000 gerações. O nome vem da
árvore, que o BEAUti batizou com o nome da primeira partição.
\end{itemize}
É recomendável rodar pelo menos mais uma réplica com outra semente (\ui{Random number seed:},
na janela do BEAST), em outra pasta, para conferir se as duas chegam ao mesmo resultado.

% =====================================================================
\section{Resultados}
% =====================================================================

\paragraph{Tracer.} Abra o Tracer e carregue o \arq{hylodidae.log} com o botão \ui{+} no
canto superior esquerdo. Na coluna \ui{ESS} (\emph{effective sample size}), procure valores
acima de 200; valores menores aparecem em destaque e indicam que a cadeia precisa rodar por
mais tempo. Se você rodou duas réplicas, carregue as duas e compare as médias.

\paragraph{TreeAnnotator.} O TreeAnnotator resume as milhares de árvores do arquivo
\arq{.trees} em uma só, a árvore de máxima credibilidade de clados (\emph{Maximum clade
credibility tree}).
\begin{enumerate}
\item Em \ui{Burn in percentage}, deixe 10 (as primeiras 10\% das árvores são descartadas).
\item Em \ui{Target tree type}, deixe \ui{Maximum clade credibility tree}.
\item Em \ui{Node heights}, troque \ui{Common Ancestor heights} por \ui{Mean heights}.
\item Em \ui{Input Tree File}, escolha \arq{hylodidae-p16S.trees} (\ui{Choose file...}); em
\ui{Output File}, escreva \arq{hylodidae_mcc.tree}. Clique em \ui{Run}.
\end{enumerate}

\paragraph{FigTree.} Abra o \arq{hylodidae_mcc.tree} no FigTree. No painel da esquerda:
\begin{itemize}
\item marque \ui{Node Labels} e, em \ui{Display}, escolha \ui{posterior}: é a probabilidade
posterior de cada clado;
\item marque \ui{Node Bars} e, em \ui{Display}, escolha \ui{height\_95\%\_HPD}: as barras
mostram o intervalo de 95\% de maior densidade da idade de cada nó;
\item marque \ui{Scale Axis} e, dentro dele, \ui{Reverse axis}, para ler as idades em milhões
de anos antes do presente.
\end{itemize}

% =====================================================================
\section*{Perguntas}
\addcontentsline{toc}{section}{Perguntas}
% =====================================================================

\begin{enumerate}
\item Todos os parâmetros chegaram a ESS acima de 200? Quais ficaram abaixo? O que você faria
para resolver?
\item Qual a idade estimada (média e intervalo de 95\%) do ancestral comum mais recente de
Hylodidae?
\item No Tracer, compare as médias de \verb|clockRate.COI| e \verb|clockRate.p16S|. Quantas
vezes o COI evolui mais rápido que o 16S? O resultado está dentro do que esperávamos pela
\emph{prior}?
\end{enumerate}

\begin{resposta}
Valores de duas corridas de um milhão de gerações (sementes 1 e 2), com 10\% de
\emph{burn-in}; a sua corrida terá números parecidos, mas não idênticos.

1. Não. Com um milhão de gerações, a maioria dos parâmetros fica abaixo de 200: a posterior
(cerca de 130), a altura da árvore (135 e 184), as taxas dos relógios (cerca de 90 a 130) e
vários parâmetros do bModelTest, os mais baixos (\verb|BMT_ProportionInvariable.p16S| e
\verb|BMT_gammaShape.p16S|, cerca de 40). Ao todo, cerca de 30 dos 44 parâmetros ficaram
abaixo de 200; entre os que passaram estão a verossimilhança (222 e 280), a taxa de especiação
(cerca de 250) e alguns parâmetros do bModelTest. A solução é rodar mais gerações (pelo menos
$10^7$) ou combinar várias réplicas com o LogCombiner.

2. Cerca de 42~Ma: 42,8~Ma (95\% HPD 22,1--69,3) e 41,3~Ma (20,4--66,6) nas duas corridas.
O intervalo é largo porque a única informação de tempo é a \emph{prior} da taxa do 16S.

3. \verb|clockRate.p16S| $\approx 0{,}0050$ e \verb|clockRate.COI| $\approx 0{,}0178$: o COI
evolui cerca de 3,5 vezes mais rápido que o 16S (95\% de 2,6 a 4,8 vezes). É mais que as 2
vezes de Vences et al., que usaram distâncias não corrigidas (que subestimam as diferenças no
gene mais rápido), mas fica bem dentro do intervalo da \emph{prior} (0,3 a 14 vezes).
\end{resposta}

% =====================================================================
\appendix
\section{Problemas comuns}\label{ap:problemas}
% =====================================================================

\begin{description}[style=nextline,leftmargin=1.5em]
\item[\ui{BEAST Model Test} não aparece] O pacote não foi instalado, ou o BEAUti não foi
reaberto depois da instalação.
\item[O BEAUti mostra uma partição só] O bloco \verb|begin sets; ... end;| não foi colado no fim
do arquivo, ou falta algum ponto e vírgula.
\item[A taxa do relógio não pode ser estimada] Desmarque \ui{Automatic set clock rate} no menu
\ui{Mode}.
\item[O BEAST avisa que o arquivo \arq{.log} já existe] Já há uma análise com o mesmo nome na
pasta. Rode a nova réplica em outra pasta, ou escolha a opção de sobrescrever.
\end{description}

% =====================================================================
\section{Alterações em relação à versão anterior}\label{ap:mudancas}
% =====================================================================

\begin{enumerate}
\item \textbf{Dados.} No lugar do 16S de \emph{Pristimantis rupicola}, é usada a matriz
concatenada de Hylodidae (16S + COI) da Prática~4, com duas partições e árvore compartilhada.
\item \textbf{Relógios.} Um relógio estrito para cada marcador: o do 16S, com a \emph{prior}
do tutorial original, calibra o tempo; o do COI é estimado, com uma \emph{prior} lognormal
larga centrada no dobro da taxa do 16S.
\item \textbf{Simplificações para a aula.} Relógio estrito no lugar do relaxado, Yule no lugar
do \emph{Birth Death} e um milhão de gerações no lugar de dez milhões, com quadros indicando o
que seria feito numa análise real.
\item \textbf{Monofilia.} Hylodidae forçado como monofilético, por haver um único grupo externo.
\item \textbf{Visualização.} Detalhes do TreeAnnotator e do FigTree (barras de 95\% HPD e eixo
de tempo); capturas de tela substituídas por descrições em texto; perguntas ao final.
\end{enumerate}

\begin{thebibliography}{9}
\bibitem{bouckaert2017} Bouckaert, R.~R. \& Drummond, A.~J. 2017. bModelTest: Bayesian
phylogenetic site model averaging and model comparison. \emph{BMC Evolutionary Biology} 17: 42.
\bibitem{bouckaert2019} Bouckaert, R. \emph{et al.} 2019. BEAST 2.5: an advanced software
platform for Bayesian evolutionary analysis. \emph{PLoS Computational Biology} 15(4): e1006650.
\bibitem{desa2022} de Sá, F.~P., Condez, T.~H., Lyra, M.~L., Haddad, C.~F.~B. \& Malagoli,
L.~R. 2022. Unveiling the diversity of Giant Neotropical Torrent frogs (Hylodidae):
phylogenetic relationships, morphology, and the description of two new species.
\emph{Systematics and Biodiversity} 20(1): 2039318.
\bibitem{dufresnes2022} Dufresnes, C. \emph{et al.} 2022. Shedding light on taxonomic chaos:
diversity and distribution of South Asian skipper frogs (Anura, Dicroglossidae,
\emph{Euphlyctis}). \emph{Systematics and Biodiversity} 20(1): 2102686.
\bibitem{fouquet2012} Fouquet, A., Noonan, B.~P., Rodrigues, M.~T., Pech, N., Gilles, A. \&
Gemmell, N.~J. 2012. Multiple Quaternary refugia in the eastern Guiana Shield revealed by
comparative phylogeography of 12 frog species. \emph{Systematic Biology} 61(3): 461--489.
\bibitem{gehara2017} Gehara, M., Barth, A., Oliveira, E.~F., Costa, M.~A., Haddad, C.~F.~B.
\& Vences, M. 2017. Model-based analyses reveal insular population diversification and cryptic
frog species in the \emph{Ischnocnema parva} complex in the Atlantic forest of Brazil.
\emph{Molecular Phylogenetics and Evolution} 112: 68--78.
\bibitem{lemmon2007} Lemmon, E.~M., Lemmon, A.~R. \& Cannatella, D.~C. 2007. Geological and
climatic forces driving speciation in the continentally distributed trilling chorus frogs
(\emph{Pseudacris}). \emph{Evolution} 61(9): 2086--2103.
\bibitem{rambaut2018} Rambaut, A., Drummond, A.~J., Xie, D., Baele, G. \& Suchard, M.~A.
2018. Posterior summarization in Bayesian phylogenetics using Tracer 1.7. \emph{Systematic
Biology} 67(5): 901--904.
\bibitem{vences2005b} Vences, M., Thomas, M., Bonett, R.~M. \& Vieites, D.~R. 2005.
Deciphering amphibian diversity through DNA barcoding: chances and challenges.
\emph{Philosophical Transactions of the Royal Society B} 360(1462): 1859--1868.
\bibitem{vittorazzi2021} Vittorazzi, S.~E., Augusto-Alves, G., Neves-da-Silva, D.,
Carvalho-e-Silva, A.~M.~P.~T., Recco-Pimentel, S.~M., Toledo, L.~F., Lourenço, L.~B. \&
Bruschi, D.~P. 2021. Paraphyly in the giant torrent-frogs (Anura: Hylodidae:
\emph{Megaelosia}) and the description of a new genus. \emph{Salamandra} 57(2): 274--284.
\end{thebibliography}

\end{document}
